\documentclass[11pt]{article}
\usepackage[utf8]{inputenc}
\usepackage{amsmath,amssymb}
\usepackage[margin=1in]{geometry}
\usepackage{hyperref}
\emergencystretch=3em

\title{Measurability of the chart integral in the parameter}
\author{Claude, with Daniel Murfet}
\date{2026-09-07}

\begin{document}
\maketitle
\sloppy

\begin{abstract}
Astra~\#10 rank~2. For a coefficient family $c:\Omega\to\mathbb N^2\to\mathbb R\to\mathbb R$ jointly
measurable in $(w,s)$ with a common radius $\rho>b$, the parametrised amplitude
$(w,z)\mapsto\Phi_w(z_1,z_2,Nz_1^{k_1}z_2^{k_2})$ is measurable: its rectangular partial sums are
measurable and converge everywhere on the box, so the limit is measurable
(\texttt{anaAmp\_param\_measurable}); the product integrand follows, and the parametric Bochner
integral over the box measure gives measurability of $w\mapsto Z_w(N)$
(\texttt{twoDAmp\_param\_measurable}), discharging the first measurability hypothesis of the
averaged Taylor tree. Zero \texttt{sorry}/\texttt{axiom}.
\end{abstract}

\section{The things formalised}
{\small
\begin{verbatim}
theorem anaAmp_param_measurable {Ω : Type*} [MeasurableSpace Ω] (b N ρ : ℝ) (k₁ k₂ : ℕ)
    (hb : 0 < b) (hbρ : b < ρ) (c : Ω → ℕ × ℕ → ℝ → ℝ)
    (hcm : ∀ ij, Measurable fun p : Ω × ℝ => c p.1 ij p.2) (H : Ω → ℝ → ℝ)
    (hc : ∀ w ij s, |c w ij s| * ρ ^ (ij.1 + ij.2) ≤ H w s) :
    Measurable fun p : Ω × (ℝ × ℝ) =>
      anaAmp (c p.1) b p.2.1 p.2.2 (N * (p.2.1 ^ k₁ * p.2.2 ^ k₂))
theorem twoDAmp_param_measurable {Ω : Type*} [MeasurableSpace Ω] (β b N ρ : ℝ)
    (h₁ h₂ k₁ k₂ : ℕ) (hb : 0 < b) (hbρ : b < ρ) (c : Ω → ℕ × ℕ → ℝ → ℝ)
    (hcc : ∀ w ij, Continuous (c w ij))
    (hcm : ∀ ij, Measurable fun p : Ω × ℝ => c p.1 ij p.2)
    (H : Ω → ℝ → ℝ) (hH : ∀ w, Continuous (H w))
    (hc : ∀ w ij s, |c w ij s| * ρ ^ (ij.1 + ij.2) ≤ H w s) :
    Measurable fun w => twoDAmp β b N h₁ h₂ k₁ k₂ (anaAmp (c w) b)
theorem twoDAmp_param_aestronglyMeasurable ... :
    AEStronglyMeasurable (fun w => twoDAmp β b N h₁ h₂ k₁ k₂ (anaAmp (c w) b)) μ
\end{verbatim}
}

\section{Mathematics}

Each term $c_w(ij,s)\,\bar u^i\bar v^j$ ($\bar u,\bar v$ the clamped coordinates, $s=Nu^{k_1}v^{k_2}$)
is measurable in $(w,z)$; the finite partial sums over $F\subset\mathbb N^2$ are measurable; for every
$(w,z)$ the series is absolutely convergent since $\bar u,\bar v\le b<\rho$ and $|c|\rho^{i+j}\le H$,
so the partial sums converge to the amplitude along the directed family of finite sets, a countably
generated filter; hence the amplitude is measurable as a pointwise limit. The integrand
$u^{h_1}v^{h_2}e^{-\beta s^2}\Phi_w$ is then measurable on $\Omega\times\mathbb R^2$, and the integral
over the finite box measure is measurable in $w$ by the Fubini-type measurability of parametric
integrals; finally $Z_w(N)$ equals that integral by the product-form identity of unit~89.

\section{Paper-facing meaning}

Ordinary but necessary: the per-stratum integral $\int Z_w(N)\rho_I(w)\,dw$ of \texttt{thm:TaylorTree}
is well defined once the chart amplitude depends measurably (in the natural joint sense) on the
noncritical coordinates, which holds for analytic $\xi,\eta$. Together with unit~137 this makes the
averaged Taylor tree a theorem about a measurable family with a common envelope, modulo the
measurability of the canonical coefficients (Astra~\#10 rank~3, next).

\section{Lean notes}

\begin{itemize}
\item \texttt{measurable\_of\_tendsto\_metrizable' atTop} works along the filter of finite subsets
of $\mathbb N^2$ because \texttt{atTop.isCountablyGenerated} applies to any countable preorder; the
\texttt{HasSum} of each point is exactly the required \texttt{Tendsto} of partial sums, packaged with
\texttt{tendsto\_pi\_nhds}.
\item \texttt{StronglyMeasurable.integral\_prod\_right'} needs \texttt{SFinite} of the second
measure; for the project's \texttt{boxMeasure₂} (a plain \texttt{def}) provide it with a local
\texttt{have ... := by unfold boxMeasure₂; infer\_instance}.
\item The linter prefers \texttt{have} to \texttt{haveI} for propositional instances.
\end{itemize}

\end{document}
